\chapter{Capa de Control y Abstracción de Hardware}
\label{ch:control-hal}

\section{Propósito de Este Capítulo}
\label{sec:control-hal-purpose}

Este capítulo documenta cómo el firmware \textbf{MICSBOL/ESP32\_BT\_Controller} separa:

\begin{itemize}
    \item \textbf{la lógica de control} (qué debería hacer el robot cuando el usuario mueve un stick), y
    \item \textbf{la conducción del hardware} (cómo el ESP32 escribe señales PWM/GPIO/I\textsuperscript{2}C).
\end{itemize}

El repositorio implementa esta separación principalmente mediante:

\begin{itemize}
    \item \texttt{control.cpp} / \texttt{control.h} --- la capa de control
    \item \texttt{hal\_output.cpp} --- la capa de abstracción de hardware (HAL)
    \item \texttt{pins.h} --- el mapeo de hardware utilizado por el HAL
    \item \texttt{pulse.cpp} / \texttt{pulse.h} --- salidas de pulso no bloqueantes usadas por acciones de evento
\end{itemize}

\section{Arquitectura en Capas (Vista de Alto Nivel)}
\label{sec:control-hal-architecture}

El firmware usa un flujo en capas desde la app Android hasta el hardware físico:

\begin{enumerate}
    \item \textbf{La app Android} transmite paquetes de estado RC y paquetes de evento por Bluetooth (SPP).
    \item \textbf{Capa receptora} (\texttt{receiver.cpp}) decodifica paquetes y actualiza:
    \begin{itemize}
        \item \texttt{rcStatePacket} para valores de control continuo
        \item \texttt{rcEventPacket} para eventos de botones de una sola vez
    \end{itemize}
    \item \textbf{Capa de control} (\texttt{control.cpp}) interpreta \texttt{rcStatePacket} y llama funciones del HAL.
    \item \textbf{HAL} (\texttt{hal\_output.cpp}) realiza las operaciones reales de GPIO/PWM/I\textsuperscript{2}C.
\end{enumerate}


\begin{figure}[H]
    \centering
    \begin{tikzpicture}[
        font=\small,
        >=Stealth,
        line/.style={->, very thick, draw=black!70},
        box/.style={
            draw,
            rounded corners=2mm,
            very thick,
            align=center,
            minimum width=12.6cm,
            minimum height=1.15cm,
            fill=gray!4
        },
        boxApp/.style={box, fill=blue!6,  draw=blue!55!black},
        boxCore/.style={box, fill=teal!6, draw=teal!55!black},
        boxHw/.style={box, fill=green!6, draw=green!55!black},
        tag/.style={
            draw=black!25,
            fill=white,
            rounded corners=1mm,
            inner sep=2pt,
            font=\scriptsize\bfseries,
            text=black!70
        }
        ]

        % Manual coordinates (no "below=of" required)
        \node[boxApp]  (android) at (0, 0.0)  {App Android\\\footnotesize Controles de UI (sticks, perillas, interruptores, botones)};
        \node[boxCore] (bt)      at (0,-1.7)  {Bluetooth Serial (SPP)};
        \node[boxCore] (rx)      at (0,-3.4)  {receiver.cpp\\\footnotesize Buffer circular + decodificación en rcStatePacket / rcEventPacket};
        \node[boxCore] (control) at (0,-5.1)  {control.cpp\\\footnotesize Política de control (mapear entradas a acciones)};
        \node[boxCore] (hal)     at (0,-6.8)  {hal\_output.cpp\\\footnotesize Salidas de hardware (PWM / GPIO / MCP)};
        \node[boxHw]   (hw)      at (0,-8.5)  {Hardware\\\footnotesize Motores, servos, LEDs, interruptores, MCP23017};

        % arrows
        \draw[line] (android) -- (bt);
        \draw[line] (bt) -- (rx);
        \draw[line] (rx) -- (control);
        \draw[line] (control) -- (hal);
        \draw[line] (hal) -- (hw);

        % left-side step tags
        \node[tag] at (-7.1, 0.0)  {1};
        \node[tag] at (-7.1,-1.7)  {2};
        \node[tag] at (-7.1,-3.4)  {3};
        \node[tag] at (-7.1,-5.1)  {4};
        \node[tag] at (-7.1,-6.8)  {5};
        \node[tag] at (-7.1,-8.5)  {6};

    \end{tikzpicture}
    \caption{Flujo de datos de control desde Android hasta el hardware del ESP32.}
\end{figure}

\section{Paquetes de Control y Estructuras de Estado}
\label{sec:control-hal-packets}

Los valores de control se almacenan en \texttt{rcStatePacket} (definido en \texttt{packets.h}).
Los campos incluyen:

\begin{itemize}
    \item \texttt{leftStickX}, \texttt{leftStickY}
    \item \texttt{rightStickX}, \texttt{rightStickY}
    \item \texttt{leftKnob}, \texttt{rightKnob}
    \item \texttt{switches} (máscara de bits)
\end{itemize}

En la lógica actual del firmware, los canales de sticks y perillas se tratan como valores de 12 bits (rango nominal 0 a 4095).

Las pulsaciones de botones se entregan mediante \texttt{rcEventPacket} (también en \texttt{packets.h}) y se manejan con
\texttt{controlHandleEvent(eventId)}.

\section{Capa de Control (\texttt{control.cpp})}
\label{sec:control-hal-control-layer}

\subsection{Responsabilidades}
\label{subsec:control-hal-control-resp}

La capa de control es responsable de:

\begin{itemize}
    \item mapear sticks y perillas a comandos de actuadores,
    \item seleccionar qué función del HAL llamar para cada salida lógica,
    \item decodificar la máscara de bits de switches a estados por salida,
    \item manejar eventos discretos de botones disparando salidas de pulso.
\end{itemize}

\subsection{Bucle principal: \texttt{controlUpdate()}}
\label{subsec:control-hal-controlUpdate}

La función principal de control es:

\begin{verbatim}
    controlUpdate()
\end{verbatim}

En la implementación del repositorio, \texttt{controlUpdate()} lee valores desde \texttt{rcStatePacket} y llama a:

\begin{itemize}
    \item \texttt{halSetSteering(...)}
    \item \texttt{halSetMotorLeft(...)}
    \item \texttt{halSetMotorRight(...)}
    \item \texttt{halSetCameraPan(...)}
    \item \texttt{halSetLed(...)}
    \item \texttt{halSetBuzzer(...)}
    \item \texttt{halSetSwitch(index, state)}
\end{itemize}

Un helper clave es la función de mapeo de servo en \texttt{control.cpp}:

\begin{verbatim}
    static uint32_t mapServo(uint16_t v) {
        return map(v, 0, 4095, 205, 410);
    }
\end{verbatim}

\subsection{Decodificación de máscara de bits de switches}
\label{subsec:control-hal-switch-mask}

Los switches se transmiten como un solo byte \texttt{rcStatePacket.data.switches}.
La capa de control convierte bits en salidas:

\begin{verbatim}
    for (int i = 0; i < 6; i++) {
        halSetSwitch(i, (sw >> i) & 1);
    }
\end{verbatim}

\subsection{Manejo de eventos: \texttt{controlHandleEvent(eventId)}}
\label{subsec:control-hal-events}

Los eventos discretos (pulsaciones de botones) se manejan mediante:

\begin{verbatim}
    controlHandleEvent(byte eventId)
\end{verbatim}

En el repositorio, los eventos se mapean a salidas de pulso mediante \texttt{halPulseSwitch(index)}. El mapeo depende del
modo de proyecto (GPIO directos vs modo de expansor MCP).

\section{Capa de Abstracción de Hardware (\texttt{hal\_output.cpp})}
\label{sec:control-hal-hal-layer}

\subsection{Responsabilidades}
\label{subsec:control-hal-hal-resp}

El HAL implementa los \emph{mecanismos} para accionar el hardware:

\begin{itemize}
    \item salida PWM mediante ESP32 LEDC usando \texttt{ledcWrite(...)}
    \item salida GPIO mediante \texttt{digitalWrite(...)}
    \item salida de expansor MCP23017 mediante \texttt{mcpWriteOutput(...)} (I\textsuperscript{2}C)
    \item pines de dirección de motor y un mapeo con banda muerta
    \item rampa de velocidad de motor usando un callback periódico de timer
    \item salidas de pulso no bloqueantes usando el módulo de pulsos
\end{itemize}

\subsection{Salidas PWM (ejemplos)}
\label{subsec:control-hal-hal-pwm}

Ejemplos de \texttt{hal\_output.cpp}:

\begin{verbatim}
    void halSetSteering(uint32_t value) { ledcWrite(PIN_PWM_STICK_LX, value); }
    void halSetLed(uint32_t value)      { ledcWrite(PIN_PWM_KNOB_L, value);  }
    void halSetBuzzer(uint32_t value)   { ledcWrite(PIN_PWM_KNOB_R, value);  }
\end{verbatim}

\subsection{Control de motor: dirección + banda muerta + rampa}
\label{subsec:control-hal-hal-motors}

El control de motor en el HAL usa:

\begin{itemize}
    \item pines de dirección \texttt{PIN\_MD\_LEFT} y \texttt{PIN\_MD\_RIGHT}
    \item pines de salida PWM \texttt{PIN\_PWM\_STICK\_LY} (izquierdo) y \texttt{PIN\_PWM\_STICK\_RY} (derecho)
\end{itemize}

El HAL interpreta valores de motor con una banda muerta alrededor del centro, y hace una rampa del PWM real del motor
hacia un valor objetivo usando una función de actualización accionada por timer.

\subsection{Salidas de switches: GPIO directos vs MCP23017}
\label{subsec:control-hal-hal-switches}

El HAL proporciona una interfaz uniforme:

\begin{verbatim}
    halSetSwitch(uint8_t index, bool state)
\end{verbatim}

pero la implementación depende del modo de compilación:

\begin{itemize}
    \item En modo GPIO directo, se escribe un pequeño arreglo de pines ESP32 usando \texttt{digitalWrite(...)}.
    \item En modo MCP, la salida se escribe usando \texttt{mcpWriteOutput(...)} sobre I\textsuperscript{2}C.
\end{itemize}

\subsection{Salidas de pulso (acciones momentáneas)}
\label{subsec:control-hal-hal-pulses}

Las salidas de pulso se disparan mediante \texttt{halPulseSwitch(index)} y se implementan usando:

\begin{itemize}
    \item \texttt{pulseStart(pin, durationMs)} para iniciar un pulso
    \item \texttt{pulseUpdate()} para terminarlo después de la duración solicitada
\end{itemize}

Este diseño no es bloqueante, por lo que el bucle de control puede seguir ejecutándose mientras los pulsos están activos.

\section{Mapeo de Pines (\texttt{pins.h}) y Portabilidad}
\label{sec:control-hal-pins}

Los mapeos de hardware están centralizados en \texttt{pins.h}. Al adaptar el firmware a nuevo hardware:

\begin{enumerate}
    \item actualiza \texttt{pins.h} para que coincida con tu cableado,
    \item ajusta \texttt{hal\_output.cpp} si tu driver de motor o la lógica de salida difiere,
    \item mantén \texttt{control.cpp} estable como ``política'' siempre que sea posible.
\end{enumerate}

\section{Resumen}
\label{sec:control-hal-summary}

El repositorio implementa una separación clara entre:

\begin{itemize}
    \item \textbf{Política:} \texttt{control.cpp} interpreta \texttt{rcStatePacket} y decide qué acciones tomar.
    \item \textbf{Mecanismo:} \texttt{hal\_output.cpp} acciona hardware PWM/GPIO/I\textsuperscript{2}C para realizar esas acciones.
\end{itemize}

Esto hace que el firmware sea más fácil de portar entre diferentes robots y esquemas de cableado, manteniendo la base de código
organizada y mantenible.